% ==============================================================
\section{重複排除の正当性の証明}\label{app:dedup_proof}
% ==============================================================

\begin{proof}[Claim~\ref{prop:dedup}の証明]
独立出現仮定のもとでは，イベント $e$ の影響下で
$P^*$ の期待サポートは $(p_{\text{base}} + \beta)^l$，
非イベント期間では $p_{\text{base}}^l$ である．
$r = (p_{\text{base}} + \beta) / p_{\text{base}} > 1$ とおくと，
$\text{magnitude}(P^*) \propto p_{\text{base}}^l (r^l - 1)$ である．
副次パターン $P'$ は $s$ 個のブーストアイテムと $l - s$ 個の非ブーストアイテムを持つため，
$\text{magnitude}(P') \propto p_{\text{base}}^l (r^s - 1)$ となる．
$r > 1$ かつ $l > s$ より $r^l > r^s$ であるから，
$\text{magnitude}(P^*) > \text{magnitude}(P')$ が任意のパターン長 $l$ で成立する．
なお，この結果は独立出現の仮定に依存しており，
アイテム間に強い相関を持つデータでは順序が逆転しうる．
\end{proof}

% ==============================================================
\section{計算量分析の詳細}\label{app:complexity}
% ==============================================================

パイプライン全体の時間計算量は以下の通りである:
\begin{itemize}
\item \textbf{Step 0}（密集区間抽出，Apriori-window）: $O(|\mathcal{F}| \cdot N)$（詳細は\cite{dsaa2025}参照）
\item \textbf{Step 1}（変化点抽出＋局所サポート計算）: $O(|\mathcal{F}| \cdot \bar{C} \cdot h \cdot \log n)$，$|\mathcal{F}|$ は共起パターン数（$|P| \geq 2$），$\bar{C}$ は平均密集区間数，$h$ はレベルウィンドウ幅，$n$ は平均パターン出現数
\item \textbf{Step 2}（変化量評価）: Step 1 に含まれる（追加走査なし）
\item \textbf{Step 3--4}（帰属＋置換検定）: $O\!\left(\sum_{P} |\mathcal{I}_P| \cdot |\mathcal{E}| \cdot B \cdot C\right)$，
$B$ は置換回数，$C$ は区間あたり平均変化点数（2点固定）
\item \textbf{Step 5}（重複排除）: $O(R \cdot k)$，$R$ は有意結果数
\end{itemize}

% ==============================================================
\section{合成データ生成の詳細}\label{app:synthetic}
% ==============================================================

\subsection{デフォルト条件}

ベース語彙は200アイテム（$|\mathcal{I}| = 200$），各アイテムの出現確率
$p_{\text{base}} = 0.03$，トランザクション数 $N = 100{,}000$ である．
1トランザクションあたりの期待アイテム数は $200 \times 0.03 = 6.0$ である．

\begin{itemize}
\item \textbf{Type A}（語彙内ブースト，3パターン）:
    パターン $(5,15)$, $(25,35)$, $(45,55)$ を使用．
    各パターンのイベント持続期間は6{,}000トランザクション（全体の6\%），
    時間軸上に等間隔に配置．
    イベント期間外にも $p_{\text{base}}$ でベースラインサポートが存在し，
    イベント期間中に確率 $\beta$ で各アイテムを追加挿入する．
\item \textbf{Type B}（イベント非対応，2パターン）:
    パターン $(65,75)$, $(85,95)$ を使用．
    ランダムな6{,}000トランザクション区間で $\beta$ のブーストを付与するが，
    いずれのイベントとも対応しない．
\item \textbf{Type C}（デコイイベント，2件）:
    パターン変動と無関係なイベント（持続期間6{,}000）をランダム配置．
\end{itemize}

\subsection{構造条件}

\begin{itemize}
\item \textbf{OVERLAP}: イベントE1$=[16{,}000,28{,}000]$, E2$=[24{,}000,36{,}000]$ が
    区間$[24{,}000,28{,}000]$で重複．時間的弁別能力を検証する．
\item \textbf{CONFOUND}: Type Bパターンの活性区間を
    イベント時間窓に意図的に重ねて配置（交絡条件）．
\item \textbf{DENSE}: Type A 6パターン + Type B 4パターン +
    デコイ4件に倍増（高密度条件）．
\item \textbf{SHORT}: イベント持続期間を1{,}600に短縮（デフォルト6{,}000の約27\%；短期イベント条件）．
\end{itemize}

\subsection{EX2 アブレーションシナリオ}

\begin{itemize}
\item \textbf{Scenario A}（prox必要）:
    パターン $(5,15)$ に2つの密集区間を設定 ---
    1つはイベント近傍$[15{,}000,25{,}000]$（因果的），
    1つは遠方$[60{,}000,68{,}000]$（偶発的）．
    イベント区間 $[15{,}000,25{,}000]$．
\item \textbf{Scenario B}（mag必要）:
    パターン $(5,15)$ が2つのイベントの影響を受ける ---
    E1はブースト $\beta = 0.5$（強），E2は $\beta = 0.1$（弱）．
\end{itemize}

\subsection{Zipf分布パラメータ}

アイテム $k$（$k = 1, \ldots, 200$）の出現確率:
$p(k) = \min(C / k^{\alpha_z},\; 0.10)$．
$C$ は中央値アイテム（$k=100$）の確率が0.03になるよう設定し，
最大確率を0.10で打ち切る．

% ==============================================================
\section{Dunnhumby データセットの詳細}\label{app:dunnhumby}
% ==============================================================

Dunnhumby ``The Complete Journey'' データセット~\cite{dunnhumby2014complete}の
概要を表~\ref{tab:dunnhumby_stats}に示す．

\begin{table}[t]
\caption{Dunnhumby データセット統計量}
\label{tab:dunnhumby_stats}
\centering
\begin{tabular}{lr}
\toprule
指標 & 値 \\
\midrule
世帯数 & 2{,}500 \\
トランザクション数 & 276{,}484 \\
ユニーク商品数 & 92{,}339 \\
トランザクションあたり平均商品数 & 9.39 \\
観測期間 & 711 日 \\
キャンペーン数 & 30 \\
\bottomrule
\end{tabular}
\end{table}

30件のキャンペーンは3タイプに分類される．
本実験では集団レベルのサポート変動を分析対象とする問題定義（定義~\ref{def:attribution}）に基づき，
TypeAの5件のみを外部イベントとして使用する（\ref{sec:exp_ex4}節参照）:
\begin{itemize}
\item \textbf{TypeA}（5件: C8, C13, C18, C26, C30）: 大規模販促．
    クーポン対象商品数 3{,}700--38{,}200件，持続期間 40--56日．
\item \textbf{TypeB}（18件）: ターゲット型施策．
    クーポン対象商品数 200--2{,}000件，持続期間 32--61日．
\item \textbf{TypeC}（7件）: ニッチ型施策．
    クーポン対象商品数 200--1{,}500件，持続期間 32--162日．
\end{itemize}

パイプラインの設定は本文表~\ref{tab:default_params}と異なり，
$W = 300$, $\theta = 5$ とした．
$W$ は276Kトランザクションに対する適切な
時間解像度として選択した．

% ==============================================================
\section{帰無条件 FDR 検証の詳細（100 seeds）}\label{app:null_fdr}
% ==============================================================

帰無条件での BH 手続き FDR 制御を 100 seeds で検証した．

\paragraph{設定}
植え込み信号なし（Type A パターン 0 件），
ランダムに配置した 5 件のデコイイベント（各持続 6,000 トランザクション），
Type B 密集パターン 4 件（いずれのイベントとも時間的に無関係），
$N = 100{,}000$，$W = 1{,}000$，$\theta = 100$，$B = 5{,}000$，$\alpha = 0.10$．

\paragraph{結果}

\begin{center}
\begin{tabular}{lr}
\toprule
指標 & 値 \\
\midrule
Seeds 数 & 100 \\
偽陽性 0 件の seeds & 93/100 \\
偽陽性 1 件の seeds & 7/100 \\
偽陽性 2 件以上の seeds & 0/100 \\
平均偽陽性帰属数 & 0.07 \\
Mean per-seed FDR & 0.070 \\
FDR $\leq \alpha = 0.10$ & \textbf{YES} \\
\bottomrule
\end{tabular}
\end{center}

Mean per-seed FDR = 0.070 $\leq$ $\alpha = 0.10$ であり，
BH 手続きが帰無条件下で FDR を制御することが 100 seeds で確認された．
偽陽性 7 件（7\%）は BH 手続きの期待誤発見率 $\alpha = 0.10$ の範囲内であり，
確率的変動として説明される．

% ==============================================================
\section{パターン長 $l$ の頑健性}\label{app:pattern_length}
% ==============================================================

EX1の合成実験はすべて2-アイテムセット（$l=2$）を用いているが，
重複排除基準（$\lceil l/2 \rceil$ 多数共有）はパターン長 $l$ に依存する．
$l=3,4$ のパターンに対しても同様に機能するかを検証した．

\paragraph{設定}
各パターン長 $l \in \{2, 3, 4\}$ で，語彙内アイテムIDを $n_{\text{items}}=200$ 超
（ID 201以上）に設定し，ベース確率をゼロとした植え込み信号3件を生成する．
$N=100{,}000$，$W=1{,}000$，$\theta=100$，20 seedで評価．

\paragraph{結果}
\begin{center}
\begin{tabular}{lccc}
\toprule
$l$ & Prec & Rec & F1 (95\% CI) \\
\midrule
2 & 1.00 & 1.00 & 1.00~[1.00, 1.00] \\
3 & 1.00 & 1.00 & 1.00~[1.00, 1.00] \\
4 & 1.00 & 1.00 & 1.00~[1.00, 1.00] \\
\bottomrule
\end{tabular}
\end{center}

全パターン長でF1 = 1.00（95\% CI = [1.00, 1.00]）を達成した．
$l=3$ の重複排除閾値 $\lceil 3/2 \rceil = 2$ および
$l=4$ の閾値 $\lceil 4/2 \rceil = 2$ は，いずれも正確に機能している．

% ==============================================================
\section{アイテム間相関への頑健性}\label{app:item_correlation}
% ==============================================================

実データではビールとつまみのような強い共起相関が存在し，
Proposition~2（重複排除の正当性）の独立出現仮定が崩れる可能性がある．
Zipf分布とアイテム間相関の両条件で20 seed評価を行った．

\paragraph{設定}
\textbf{zipf\_only}: Zipf-1.0分布（独立出現）．
\textbf{zipf\_corr}: Zipf-1.0 + アイテムペア5組に相関確率0.5--0.7を付与．
$N=100{,}000$，$W=1{,}000$，$\theta=100$，20 seedで評価．

\paragraph{結果}
\begin{center}
\begin{tabular}{lccccc}
\toprule
条件 & Prec & Rec & F1 & 95\% CI & FAR \\
\midrule
zipf\_only & 0.74 & 0.92 & 0.81 & [0.70, 0.92] & 0.25 \\
zipf\_corr & 0.41 & 0.85 & 0.55 & [0.45, 0.65] & 0.20 \\
\bottomrule
\end{tabular}
\end{center}

zipf\_onlyでF1 = 0.81，zipf\_corrでF1 = 0.55と低下する．
精度低下の主因は，相関アイテムペアが生み出す追加の密集パターンが
偽陽性帰属を増やすためである．
これはProposition~2の独立出現仮定が実データで緩和される可能性を裏付けており，
アイテム間相関が強い場合のFAR = 0.20は許容可能な範囲に留まる．

% ==============================================================
\section{OVERLAP条件における交差シグナルの影響}\label{app:overlap_analysis}
% ==============================================================

OVERLAP条件（§\ref{sec:exp_ex1}）でF1分散が大きい（CI幅 = 0.44）理由を説明する．
E1 $= [16{,}000, 28{,}000]$ とE2 $= [24{,}000, 36{,}000]$ が重複する区間
$[24{,}000, 28{,}000]$ では，E1のアイテム集合とE2のアイテム集合が
\emph{同一ウィンドウ内で同時にブースト}される．
各アイテムの出現確率が $\beta = 0.3$ の場合，
交差パターン（例: E1アイテムとE2アイテムの組み合わせ）の
期待サポートは $0.3 \times 0.3 \times W = 90 \approx \theta$  となり，
統計的変動によって $\theta$ を超える短い密集区間が生じる．
重複排除（Step~5）はこれらを強い planted パターンに統合することで大半を除去するが，
一部のseedでは交差パターンが偽陽性として残存するため，精度に分散が生じる．

% ==============================================================
\section{EX4 定性ケーススタディ}\label{app:ex4_casestudy}
% ==============================================================

各キャンペーンで帰属スコア上位のパターンを表~\ref{tab:ex4_top_patterns}に示す．
C8（春季生鮮フェア）では TROPICAL FRUIT + SOFT DRINKS や TOMATOES + BERRIES が帰属された．
これらは青果・乳製品カテゴリのクーポン対象商品と整合しており，
季節性の高い購買パターンがキャンペーンと同期して上昇したことを示す．
C18（秋季特売）では BEEF + TROPICAL FRUIT と FLUID MILK + PIES が帰属され，
秋季の家庭内調理需要の高まりと整合する解釈が可能である．
C26・C30では帰属パターンにクーポン対象商品が直接含まれていないが，
クロスセル効果や間接的購買変化を反映している可能性がある．

\begin{table}[t]
\caption{EX4 キャンペーン別トップ帰属パターン（上位2件）}
\label{tab:ex4_top_patterns}
\centering
\small
\begin{tabular}{llcc}
\toprule
キャンペーン & パターン（商品カテゴリ） & スコア & クーポン整合 \\
\midrule
C8  & TROPICAL FRUIT + SOFT DRINKS & 1.710 & \checkmark \\
    & TOMATOES + BERRIES & 0.326 & \checkmark \\
\midrule
C13 & FLUID MILK + TROPICAL FRUIT & 1.659 & \checkmark \\
\midrule
C18 & BEEF + TROPICAL FRUIT & 0.878 & \checkmark \\
    & FLUID MILK + PIES & 0.793 & \checkmark \\
\midrule
C26 & FLUID MILK + TROPICAL FRUIT & 3.180 & \texttimes \\
    & BAKED BREAD + FLUID MILK & 1.134 & \texttimes \\
\midrule
C30 & TROPICAL FRUIT + FLUID MILK & 1.304 & \texttimes \\
    & FLUID MILK + BEEF & 0.921 & \texttimes \\
\bottomrule
\end{tabular}
\end{table}

% ==============================================================
\section{EX4 パラメータ感度分析とスケーラビリティ}\label{app:ex4_sensitivity}
% ==============================================================

\paragraph{$W$/$\theta$ 感度分析}
$W$ と $\theta$ の設定変化に対するロバスト性を確認するため，
表~\ref{tab:ex4_sensitivity}に6設定の結果を示す．
パターン数は $W$ と $\theta$ の増減に対して単調に変化するが，
クーポン整合率は全設定で33--80\%（中央値59\%）であり，
ランダムベースライン（34\%）を一貫して上回った．

\begin{table}[t]
\caption{EX4 $W$/$\theta$ 感度分析（Dunnhumby, $N=276{,}484$, $|E|=5$）}
\label{tab:ex4_sensitivity}
\centering
\small
\begin{tabular}{lrrrrrr}
\toprule
設定 & $W$ & $\theta$ & $|F|$ & 帰属数 & 整合率 & 時間(s) \\
\midrule
デフォルト   & 300 &  5 &    896 & 15 & 47\% &  2.7 \\
W100-$\theta$3 & 100 &  3 &  1{,}796 & 20 & 60\% &  9.5 \\
W100-$\theta$5 & 100 &  5 &    123 &  5 & 40\% &  1.3 \\
W300-$\theta$3 & 300 &  3 & 11{,}051 & 59 & 59\% & 40.6 \\
W500-$\theta$5 & 500 &  5 &  2{,}083 & 30 & 33\% &  5.5 \\
W500-$\theta$10 & 500 & 10 &    127 & 10 & 80\% &  1.3 \\
\bottomrule
\end{tabular}
\end{table}

\paragraph{スケーラビリティ}
表~\ref{tab:ex4_sensitivity} は $|F|$（パターン数）と処理時間の関係も示している．
Step~0（Apriori-window）の計算量は $O(|F| \cdot N)$ に比例し，
$|F| = 123$ で 1.3 秒，$|F| = 11{,}051$ で 38.9 秒と，
パターン数に対して概ね線形にスケールする．
Steps~1--5（帰属パイプライン）の計算量は $O(|F| \cdot |E| \cdot B \cdot N)$ であり，
$B = 5{,}000$ の置換検定がボトルネックになりうるが，
パターン・区間単位でのRayon 並列化により，
最大設定（$|F|=11{,}051$，$|E|=5$）でも 1.7 秒に収まった．

% ==============================================================
\section{パラメータ選択の根拠}\label{app:params}
% ==============================================================

主要パラメータの選択根拠を以下に述べる:

\begin{itemize}
\item \textbf{ウィンドウサイズ $W = 1{,}000$}:
    $N = 100{,}000$ トランザクションに対して1\%の時間解像度に相当する．
    $W$ が小さすぎるとサポート推定が不安定になり，
    大きすぎると変化点の時間的精度が低下する．
\item \textbf{最小サポート $\theta = 100$}:
    期待ベースラインサポート
    $W \cdot p_{\text{base}}^2 = 1{,}000 \times 0.03^2 = 0.9$ に対して十分に大きく，
    偶然の共起を除外しつつ，ブースト時のサポート増加を捕捉できる閾値．
    $\theta = W / 10 = 100$ はZipf分布下での偽陽性パターン数を抑制するための設定である．
\item \textbf{減衰パラメータ $\sigma = W$}:
    1ウィンドウ幅離れた変化点でproxが $e^{-1} \approx 0.37$ に減衰する設定．
    変化点とイベント境界の距離が $W$ を超えると帰属スコアが急速に低下する．
\item \textbf{置換回数 $B = 5{,}000$}:
    $p$ 値の解像度 $1/(B+1) \approx 0.0002$．
    BH補正により$\alpha = 0.10$で数百の仮説を同時検定する場合にも
    十分な $p$ 値精度を保証する．
\item \textbf{有意水準 $\alpha = 0.10$}:
    FDR制御水準として標準的な値．
    $\alpha = 0.05$ と比較して検出力を優先する設定であり，
    偽発見率10\%以下を許容する．
\item \textbf{変化量正規化（magnitude normalization）}:
    平方根正規化（式~(\ref{eq:mag_sqrt})）を採用する．
    絶対変化量はベースラインが高いパターンほどスコアが膨張するバイアスがあり，
    完全正規化（$\div s^{-}$）は低ベースラインパターンのスコアを過剰に増幅する．
    平方根正規化はこの両極端の中間に位置し，
    Zipf的な頻度分布を持つ実世界データに対して安定した精度を達成する．
\end{itemize}
